\clearpage
\VerdierPhysicalPageAnchor{0149}%
\begingroup
\small
\setlength{\parskip}{0pt}%
\setlength{\emergencystretch}{1.5em}%
\setlength{\abovedisplayskip}{3pt plus 1pt minus 1pt}%
\setlength{\belowdisplayskip}{3pt plus 1pt minus 1pt}%

\VerdierSemanticAnchor{ch2-2-3-5-injectivity-p149-continuation}
\noindent such that $uf=0$. Using a), i), we then see that there exists a
commutative diagram:

\VerdierSemanticAnchor{ch2-2-3-5-injectivity-refinement-diagram}
\begin{center}
  \begin{tikzcd}[column sep=3.7em,row sep=2.3em]
    && {} \\
    & {} \arrow[ur,"v"] && {t'\in S_{\mathcal A}(\mathcal A\cap\mathcal B)} \\
    X \arrow[ur,"f"] && Y \arrow[ul,"s"] \arrow[uu,"{t'}"'] &
  \end{tikzcd}
\end{center}

\VerdierSemanticAnchor{ch2-2-3-5-injectivity-conclusion}
\noindent such that $vf=0$. It follows that
$m=Q_1(t')^{-1}Q_1(vf)=0$.

\VerdierSemanticAnchor{ch2-2-3-5-surjectivity-map}
Let us now show that the homomorphism
\[
  \operatorname{Hom}_{\mathcal A/(\mathcal A\cap\mathcal B)}(X,Y)
  \longrightarrow
  \operatorname{Hom}_{\mathcal D/\mathcal B}(X,Y)
\]
is surjective.

\VerdierSemanticAnchor{ch2-2-3-5-surjectivity-roof}
An element $p\in\operatorname{Hom}_{\mathcal D/\mathcal B}(X,Y)$ may be
written in the form
\[
  p=Q_2(t)^{-1}Q_2(g),\qquad t\in S_{\mathcal D}(\mathcal B).
\]

\VerdierSemanticAnchor{ch2-2-3-5-surjectivity-refinement}
Using a), i), we see that such an element may always be written in the form
\[
  p=Q_2(s)^{-1}Q_2(g'),\qquad
  s\in S_{\mathcal A}(\mathcal A\cap\mathcal B),
\]

\VerdierSemanticAnchor{ch2-2-3-5-surjectivity-conclusion}
\noindent and hence comes from an element of
$\operatorname{Hom}_{\mathcal A/(\mathcal A\cap\mathcal B)}(X,Y)$.

\VerdierSemanticAnchor{ch2-2-3-5-c-proof-i-implies-ii}
It remains to prove the equivalence in c). Let us first prove i)
$\Rightarrow$ ii).

\VerdierSemanticAnchor{ch2-2-3-5-iterated-quotient}
By (2.3.1), the category $\mathcal D/\mathcal B$ is canonically isomorphic
to the category
$(\mathcal D/(\mathcal A\cap\mathcal B))
/(\mathcal B/(\mathcal A\cap\mathcal B))$, and the canonical functor

\VerdierSemanticAnchor{ch2-2-3-5-q-functor}
\[
  Q\colon\mathcal D/(\mathcal A\cap\mathcal B)
  \longrightarrow\mathcal D/\mathcal B
\]
is the quotient functor.

\VerdierSemanticAnchor{ch2-2-3-5-q-free-reduction}
Hence, by Proposition (2.3.3, a)), to prove the implication i)
$\Rightarrow$ ii), it suffices to show that for every object $X$ of
$\mathcal A$ and every object $B$ of $\mathcal B$,
$\operatorname{Hom}_{\mathcal D/(\mathcal A\cap\mathcal B)}(B,X)=0$.

\VerdierSemanticAnchor{ch2-2-3-5-m-representation}
An element $m$ of
$\operatorname{Hom}_{\mathcal D/(\mathcal A\cap\mathcal B)}(B,X)$ may be
written in the form $Q_3(f)Q_3(s)^{-1}$, where

\VerdierSemanticAnchor{ch2-2-3-5-q3-functor}
\[
  Q_3\colon\mathcal D\longrightarrow
  \mathcal D/(\mathcal A\cap\mathcal B)
\]
is the canonical functor, and where $s$ is an element of
$S_{\mathcal D}(\mathcal A\cap\mathcal B)$, the multiplicative system of
$\mathcal D$ associated with $\mathcal A\cap\mathcal B$ (2.2.2). We then
have a diagram in $\mathcal D$:

\VerdierSemanticAnchor{ch2-2-3-5-w-diagram}
\begin{center}
  \begin{tikzcd}[column sep=4em,row sep=2.4em]
    & W \arrow[dl,"s"'] \arrow[dr,"f"] & \\
    B && X
  \end{tikzcd}
  \quad ,
\end{center}
\endgroup
